\section{Counting representations of one}\label{sec:counting}

We prove Corollary~\ref{cor:counting} using the uniform upper bound
in Theorem~\ref{thm:main}, together with the prime number theorem and
elementary identities. Divisor-rich denominators and divisor-indexed
splitting are also used by Konyagin and Elsholtz
\cite{Konyagin2014,Elsholtz2016}. The first step is
to obtain a short distinct expansion of one with a denominator having
many divisors. Lemma~\ref{lem:distinct} concerns totals below one; at
total one we need a different repetition-removal argument, preserving
an odd divisor of a denominator.

\begin{lemma}\label{lem:marked-cleanup}
Let $Q>1$ be odd. Suppose that a finite list of positive integers
$m_1,\ldots,m_t$ satisfies $\sum_{i=1}^t1/m_i=1$ and contains a
multiple of $Q$. Then $1$ has a representation by at most $t$ distinct
positive unit fractions, still with a denominator divisible by $Q$.
\end{lemma}

\begin{proof}
Whenever a denominator $m$ occurs twice, replace those two terms by
\[
 \frac2m=\frac1{m/2}\quad(m\text{ even}),
 \qquad
 \frac2m=\frac1{(m+1)/2}+\frac1{m(m+1)/2}
       \quad(m\text{ odd}).
\]
The sum stays equal to one and the length does not increase.
There remains a denominator divisible by $Q$: if no such denominator
is removed it persists, while if $Q\mid m$, the even replacement
$m/2$ is a multiple of $Q$ because $Q$ is odd, and the larger odd
replacement is a multiple of $m$.

This invariant excludes a denominator one, since that would force the
singleton list $(1)$. It also excludes a repeated denominator two,
since two halves exhaust the sum and neither denominator is divisible
by $Q$. Thus an odd repeated denominator is at least three. In an odd
replacement the smaller new denominator is strictly below $m$ and the
larger is strictly above $m$. The sorted denominator tuple therefore
strictly decreases lexicographically: all entries below $(m+1)/2$
are unchanged, and one extra copy of $(m+1)/2$ is inserted.
At a fixed length such a descent in positive integer tuples terminates.
Indeed, the first coordinate can decrease only finitely often, and
after it stabilizes the same argument applies successively to the other
coordinates. Even replacements strictly reduce length, so only finitely
many occur. The procedure consequently terminates with distinct
denominators and preserves the required multiple of $Q$.
\end{proof}

\begin{proof}[Proof of Corollary~\ref{cor:counting}]
\emph{A short initial expansion.}
Let $r$ be a sufficiently large integer and let $Q$ be the product of
the first $r$ odd primes. The prime number theorem, already used in
Section~\ref{sec:residues}, gives
\[
 \log Q=O(r\log r),\qquad \log\log Q=O(\log r).
\]
Apply Theorem~\ref{thm:main} to $(Q-1)/Q$ and append $1/Q$.
Lemma~\ref{lem:marked-cleanup} produces a distinct expansion of one
whose denominator set $S$ has size $s=O(\log r)$ and contains a
multiple $n$ of $Q$. Fix this one set $S$ and this one $n$.
Writing $\tau(n)$ for the number of positive divisors of $n$, we have
$\tau(n)\ge 2^r$.

\emph{Branching over divisors.}
For every positive proper divisor $d$ of $n$, the identity
\begin{equation}\label{eq:divisor-split}
 \frac1n=\frac1{n+d}+\frac1{n+n^2/d}
\end{equation}
gives two distinct integer denominators, with
\[
 n<n+d<2n<n+n^2/d.
\]
For each unchanged denominator in $S\setminus\{n\}$, at most one
choice of $d$ makes it equal to $n+d$, and at most one makes it equal
to $n+n^2/d$. Thus at least
\[
 \tau(n)-1-2(s-1)\ge 2^r-2s+1
\]
choices give distinct-denominator expansions of one, all of the same
length $\ell=s+1$. They give different expansions: from any resulting
denominator set, remove the fixed set $S\setminus\{n\}$; the smaller
of the two remaining denominators is $n+d$, which recovers $d$.
Sorting each set therefore gives a different tuple counted by $F(\ell)$.

\emph{Every sufficiently large length.}
For any distinct expansion of one with at least two terms, its largest
denominator $v$ exceeds one. Replace $v$ by $v+1$ and $v(v+1)$, using
\begin{equation}\label{eq:padding-split}
 \frac1v=\frac1{v+1}+\frac1{v(v+1)}.
\end{equation}
These are the two largest denominators of the new expansion. The
operation is injective: delete those two denominators and reinsert
the second largest minus one. Iterating a fixed number of times is
therefore injective. Our initial set contains a multiple of $Q>1$,
so it is not the singleton $\{1\}$; the operation applies to all
the expansions just constructed.

Choose an absolute constant $B\ge1$ such that $\ell\le B\log r$ for
all sufficiently large $r$. For each sufficiently large integer $k$,
take $r=\lfloor\exp(k/(2B))\rfloor$. Then $\ell\le k$, and applying
\eqref{eq:padding-split} exactly $k-\ell$ times to each of our
common-length expansions gives
\[
 F(k)\ge 2^r-2s+1\ge 2^{r-1}.
\]
Since $r\ge\tfrac12\exp(k/(2B))$ for large $k$, this proves
$\log\log F(k)\ge k/(2B)+O(1)$, hence the desired lower bound
for every sufficiently large integer $k$.

\emph{The upper bound.}
Fix a tuple counted by $F(k)$ and put $P_0=1$ and
$P_i=n_1\cdots n_i$. Before the $i$th term, the positive remainder
\[
 R_i=1-\sum_{j<i}\frac1{n_j}=\sum_{j=i}^k\frac1{n_j}
\]
is a rational number with a positive integer numerator over
$P_{i-1}$. Ordering therefore gives
\[
 \frac1{P_{i-1}}\le R_i\le\frac{k}{n_i},
 \qquad n_i\le kP_{i-1},\qquad P_i\le kP_{i-1}^2.
\]
Inductively $P_i\le k^{2^i-1}$ and $n_i\le k^{2^{i-1}}$.
Counting all integer choices in this larger box shows that $F(k)$ is
finite and
\[
 F(k)\le\prod_{i=1}^k k^{2^{i-1}}=k^{2^k-1}.
\]
Together with the lower bound, this yields
$\log\log F(k)\le k\log2+\log\log k=O(k)$ for sufficiently
large $k$, as required.
\end{proof}
